% !TEX TS-program = pdflatex
% CV ciblé — Sécurité SAP / SAP BASIS
\documentclass[a4paper,10pt]{article}
\usepackage[utf8]{inputenc}
\usepackage[T1]{fontenc}
\usepackage[scaled=0.94]{helvet}
\renewcommand{\familydefault}{\sfdefault}
\usepackage[left=1.2cm,right=1.2cm,top=0.7cm,bottom=0.6cm]{geometry}
\usepackage{enumitem}
\usepackage{titlesec}
\usepackage{xcolor}
\usepackage[hidelinks]{hyperref}
\definecolor{navy}{RGB}{23,54,93}
\pagestyle{empty}
\setlength{\parindent}{0pt}
\setlength{\parskip}{0pt}
\linespread{0.82}
\hypersetup{pdftitle={CV - Baha Eddine Belhaj Mouldi - Ingénieur Sécurité SAP},pdfauthor={Baha Eddine Belhaj Mouldi},pdfsubject={Sécurité SAP, SAP BASIS, Détection de menaces, Fiori}}
\titleformat{\section}{\normalsize\bfseries\color{navy}}{}{0pt}{\MakeUppercase}[{\vspace{1pt}\color{navy}\hrule height 0.6pt}]
\titlespacing*{\section}{0pt}{3pt}{1.5pt}
\setlist[itemize]{leftmargin=1.1em,label=\textbullet,itemsep=0.1pt,topsep=0.2pt,parsep=0pt}
\newcommand{\cventry}[4]{\textbf{#1} \hfill \textbf{#4}\\[1pt]#2{\ifx\relax#3\relax\else, #3\fi}\par\vspace{1pt}}
\begin{document}
\begin{center}
{\LARGE\bfseries\color{navy} Baha Eddine Belhaj Mouldi}\\[3pt]
{\large Ingénieur SAP BASIS \& Sécurité}\\[5pt]
Tunis, Tunisie \quad|\quad +216 55 128 982 \quad|\quad \href{mailto:bahaeddine.belhajmouldi@gmail.com}{bahaeddine.belhajmouldi@gmail.com}\\[1pt]
\href{https://www.linkedin.com/in/bahamouldi}{linkedin.com/in/bahamouldi} \quad|\quad \href{https://github.com/bahamouldi}{github.com/bahamouldi} \quad|\quad \href{https://bahamouldi.github.io/portfolio/}{bahamouldi.github.io/portfolio}
\end{center}
\vspace{1pt}
\section{Profil}
Ingénieur SAP BASIS et sécurité avec expérience pratique chez Streamlink (2 mois de stage). Création d'un système automatisé de détection de menaces SAP (Python/RFC SDK + ELK + N8N + TheHive), gestion des utilisateurs/rôles (SU01, PFCG), monitoring système (SM50-SM58, ST02/ST03N), déploiement d'applications Fiori (SICF/OData), transports (TMS/STMS), sauvegardes (DB13) et gestion des correctifs (SUM). Connaissances SAP S/4HANA, SAP BTP, méthodologies ASAP/Activate et conformité (SOX, ISO 27001).
\section{Compétences clés}
\textbf{Administration SAP} : SU01 (utilisateurs), PFCG (rôles), analyse SoD, profils d'autorisation, gestion des mandants, RZ10/RZ11\\[0.5pt]
\textbf{Monitoring SAP} : SM50/SM51 (processus), SM04 (sessions), SM21 (logs), SM37/SM36 (jobs), SM12 (verrous), SM58 (RFC), ST02 (buffers), ST03N (charge)\\[0.5pt]
\textbf{Sécurité \& Conformité} : Détection transactions sensibles (SU01, SE16N, PFCG), scoring de risque, blocage auto via BAPI, SOX/ISO 27001/RGPD\\[0.5pt]
\textbf{Fiori \& Transports} : SICF (services ICF), activation OData (/n/iwfnd/maint\_service), Fiori Apps Library, TMS/STMS\\[0.5pt]
\textbf{Stack sécurité} : Python/RFC SDK, ELK Stack (Elasticsearch, Kibana, Metricbeat), N8N (orchestration), TheHive (gestion d'incidents)\\[0.5pt]
\textbf{Paysage SAP} : Environnements DEV/QAS/PRD, méthodologies ASAP \& SAP Activate, SAP S/4HANA, SAP BTP, sauvegardes DB13/BRTools, correctifs SUM\\[0.5pt]
\textbf{Outils \& Systèmes} : Wireshark, Nmap, Burp Suite, Linux/Kali, Bash, PowerShell, Docker, Git
\section{Expérience professionnelle}
\cventry{Stagiaire Cybersécurité — Détection de menaces SAP}{Streamlink}{Tunisie}{07/2025 -- 08/2025}
\begin{itemize}
  \item Système automatisé : scripts Python + RFC SDK pour collecte temps réel, ELK pour centralisation, N8N pour orchestration des alertes, TheHive pour suivi des incidents. Détection $<$2 min, blocage auto $<$7 min, 0\% faux positifs.
  \item Gestion utilisateurs/rôles (SU01, PFCG) : rôle ``Monitoring Basis'' avec 10 transactions (SM50-SM58, ST03N, AL08, DB02), moindre privilège et conformité SoD.
  \item Surveillance transactions sensibles (SU01, SE16N, PFCG, SM59, SE38) : catalogue de 7 tcodes critiques, scoring automatisé et alerting temps réel via dashboards Kibana.
  \item Déploiement d'applications Fiori : activation services ICF (SICF), services OData, création de rôles PFCG avec Business Catalogs/Groups/Spaces pour Fiori Launchpad.
  \item Transports via TMS/STMS (DEV $\rightarrow$ QAS $\rightarrow$ PRD), planification sauvegardes (DB13), analyse correctifs (SUM) et tuning profils système (RZ10/RZ11, ST02).
\end{itemize}
\cventry{Spécialiste Cybersécurité Junior --- Stagiaire PFE}{Digital Power Consulting}{Tunisie}{01/2026 -- 05/2026}
\begin{itemize}
  \item Tests d'intrusion sur 8 applications web et API, production de plusieurs rapports de remédiation détaillés.
\end{itemize}
\cventry{Stagiaire Sécurité réseau}{Tunisie Telecom}{Tunisie}{07/2024}
\begin{itemize}
  \item Évaluation sécurité réseau, identification de 10+ vulnérabilités, projet détection SOC avec Elastic/Kibana et règles Sigma.
\end{itemize}
\section{Projets clés}
\cventry{Système de détection et réponse aux menaces SAP}{Python, PyRFC, ELK, N8N, TheHive}{}{2025}
\begin{itemize}
  \item Pipeline SOC complet : Python/RFC SDK collecte les logs des tables SAP (USR02, UST04, STOR) $\rightarrow$ Metricbeat $\rightarrow$ Elasticsearch $\rightarrow$ dashboards Kibana (sécurité, performance, activité) $\rightarrow$ workflows N8N $\rightarrow$ cas TheHive.
  \item Surveillance de transactions sensibles : catalogue de 7 tcodes critiques (SU01, SE16N, PFCG, SM59, SE38, SA38, SE10) avec niveaux de risque (CRITIQUE/ÉLEVÉ/MOYEN), scoring automatisé par activité utilisateur.
  \item Réponse auto : BAPI\_USER\_EXISTENCE\_CHECK + BAPI\_USER\_LOCK + BAPI\_TRANSACTION\_COMMIT pour blocage instantané. 3 scénarios testés : connexions anormales, escalade de privilèges, accès non autorisé à SU01.
  \item Résultats : détection $<$2 min, cycle de réponse $<$7 min, 0\% faux positifs, couverture 100\%. Panels Kibana : connexions échouées/heure, heatmap transactions sensibles, distribution risques, alertes temps réel.
\end{itemize}
\cventry{BeeWAF --- Pare-feu applicatif intelligent (PFE)}{FastAPI, ML, Docker, Kubernetes, ELK}{}{2026}
\begin{itemize}
  \item WAF d'entreprise : 10 041 règles + 3 modèles ML (RandomForest, GradientBoosting, IsolationForest). 27 modules : anti-bot, DLP, anti-DDoS, virtual patching (37 CVE).
  \item Conformité 7 référentiels (OWASP, PCI DSS, GDPR, NIST, ISO 27001, CIS, SOC 2). Score 98,2/100, 0\% faux positifs sur tous les jeux de test.
  \item Architecture : backend FastAPI, déploiement Docker/Kubernetes, dashboards de monitoring ELK, alerting automatisé pour équipes SOC d'entreprise.
\end{itemize}
\section{Formation}
\cventry{Diplôme d'Ingénieur en Génie Informatique}{ENIT}{Tunisie}{09/2023 -- 06/2026}
\cventry{Classes préparatoires (Physique-Technologie)}{IPEI El Manar}{Tunisie}{09/2021 -- 06/2023}
\section{Certifications}
Réseaux --- Cisco (2025) ; Deep Learning --- NVIDIA (2025) ; Adversarial ML --- NVIDIA (2025) ; Machine Learning --- NVIDIA (2024) ; Git \& GitHub --- 365 Data Science (2024).
\section{Langues}
Anglais (B2) \quad|\quad Français (B2) \quad|\quad Arabe (Natif) \quad|\quad Chinois (Notions)
\end{document}
